\documentclass[10pt,a4paper]{amsart}
\usepackage{amsmath,amssymb,amsthm,graphicx}
% Retained minimal preamble for the Stacks Project chapter more-algebra.tex (More on Algebra).
% The chapter file pinned-original-source.tex carries only \input{preamble} on its first line, so no
% preamble text lives inside the chapter itself.  The Stacks Project uses one single shared preamble for
% every chapter; the snapshot stages it identically as sources/stacks-chow/preamble.tex and
% sources/stacks-etale-cohomology/preamble.tex.  The block below is copied mechanically, line for line,
% from sources/stacks-chow/preamble.tex lines 14-18 (the Xy-pic setup) and 154-217 (the native theorem
% declarations together with the mathematical macros).  Nothing was added, removed or redefined; the
% wrapper class and the amsmath/amssymb/amsthm/graphicx base packages are supplied by the assembler.
% Chapter: more-algebra.tex -- More on Algebra.
% For commutative diagrams we use Xy-pic
\usepackage[all]{xy}

% We use 2cell for 2-commutative diagrams.
\xyoption{2cell}

% Theorem environments.
%
\theoremstyle{plain}
\newtheorem{theorem}[subsection]{Theorem}
\newtheorem{proposition}[subsection]{Proposition}
\newtheorem{lemma}[subsection]{Lemma}

\theoremstyle{definition}
\newtheorem{definition}[subsection]{Definition}
\newtheorem{example}[subsection]{Example}
\newtheorem{exercise}[subsection]{Exercise}
\newtheorem{situation}[subsection]{Situation}

\theoremstyle{remark}
\newtheorem{remark}[subsection]{Remark}
\newtheorem{remarks}[subsection]{Remarks}

\numberwithin{equation}{subsection}

% Macros
%
\def\lim{\mathop{\mathrm{lim}}\nolimits}
\def\colim{\mathop{\mathrm{colim}}\nolimits}
\def\Spec{\mathop{\mathrm{Spec}}}
\def\Hom{\mathop{\mathrm{Hom}}\nolimits}
\def\Ext{\mathop{\mathrm{Ext}}\nolimits}
\def\SheafHom{\mathop{\mathcal{H}\!\mathit{om}}\nolimits}
\def\SheafExt{\mathop{\mathcal{E}\!\mathit{xt}}\nolimits}
\def\Sch{\mathit{Sch}}
\def\Mor{\mathop{\mathrm{Mor}}\nolimits}
\def\Ob{\mathop{\mathrm{Ob}}\nolimits}
\def\Sh{\mathop{\mathit{Sh}}\nolimits}
\def\NL{\mathop{N\!L}\nolimits}
\def\CH{\mathop{\mathrm{CH}}\nolimits}
\def\proetale{{pro\text{-}\acute{e}tale}}
\def\etale{{\acute{e}tale}}
\def\QCoh{\mathit{QCoh}}
\def\Ker{\mathop{\mathrm{Ker}}}
\def\Im{\mathop{\mathrm{Im}}}
\def\Coker{\mathop{\mathrm{Coker}}}
\def\Coim{\mathop{\mathrm{Coim}}}

% Boxtimes
%
\DeclareMathSymbol{\boxtimes}{\mathbin}{AMSa}{"02}

%
% Macros for moduli stacks/spaces
%
\def\QCohstack{\mathcal{QC}\!\mathit{oh}}
\def\Cohstack{\mathcal{C}\!\mathit{oh}}
\def\Spacesstack{\mathcal{S}\!\mathit{paces}}
\def\Quotfunctor{\mathrm{Quot}}
\def\Hilbfunctor{\mathrm{Hilb}}
\def\Curvesstack{\mathcal{C}\!\mathit{urves}}
\def\Polarizedstack{\mathcal{P}\!\mathit{olarized}}
\def\Complexesstack{\mathcal{C}\!\mathit{omplexes}}
% \Pic is the operator that assigns to X its picard group, usage \Pic(X)
% \Picardstack_{X/B} denotes the Picard stack of X over B
% \Picardfunctor_{X/B} denotes the Picard functor of X over B
\def\Pic{\mathop{\mathrm{Pic}}\nolimits}
\def\Picardstack{\mathcal{P}\!\mathit{ic}}
\def\Picardfunctor{\mathrm{Pic}}
\def\Deformationcategory{\mathcal{D}\!\mathit{ef}}

% --- appended 2026-10-12 (Chief): the project own definitions of the reference/slogan/history
% display environments, copied mechanically from the shared Stacks preamble lines 8-12
% (\usepackage{verbatim} plus the three \newenvironment lines, which the chapter preamble keeps
% verbatim from the project preamble). Nothing added, removed or redefined: these environments
% carry no mathematical content and the project itself typesets them as comments. Without these
% three lines a counted statement whose own text uses them fails with Environment slogan undefined.
\usepackage{verbatim}
\newenvironment{reference}{\comment}{\endcomment}
\newenvironment{slogan}{\comment}{\endcomment}
\newenvironment{history}{\comment}{\endcomment}
\usepackage{cleveref}
\crefname{originallemma}{original lemma}{original lemmas}
\Crefname{originallemma}{Original lemma}{Original lemmas}
\makeatletter
\newlabel{lemma-localize-bezout}{{37834}{source}{}{sourcepointer-1}{}}
\newlabel{lemma-localize-bezout@cref}{{[originallemma][0][]37834}{[1][0][]source}{}{}{}}
\makeatother
\crefname{originallemma}{original lemma}{original lemmas}
\Crefname{originallemma}{Original lemma}{Original lemmas}
\makeatletter
\newlabel{lemma-warfield}{{37737}{source}{}{sourcepointer-2}{}}
\newlabel{lemma-warfield@cref}{{[originallemma][0][]37737}{[1][0][]source}{}{}{}}
\makeatother
\title{Over a Bezout domain every finite submodule of a free module is free and every finitely presented module splits as free plus torsion}
\author{The Stacks Project Authors}
\date{stacks-project:more-algebra.tex:local-snapshot-20261012; GNU FDL 1.2 or later}
\begin{document}
\maketitle
\setcounter{subsection}{1}
\noindent\textit{Editorial note.} All delivered mathematical text is the author's own, copied from the pinned Stacks Project chapter file more-algebra.tex; no statement or proof step was rewritten, reordered or removed.\par\medskip
\noindent\textit{Editorial note.} The chapter file pinned-original-source.tex carries only \textbackslash{}input\{preamble\} on its first line, so no preamble text lives inside the chapter. The retained file sources/stacks-more-algebra/retained-preamble.tex is a mechanical, line-for-line copy of the project's single shared preamble as staged in the snapshot; nothing was added, removed or redefined and no package was added.\par\medskip
\noindent\textit{Editorial note.} The source URL is the source-level Stacks Project URL, as required when more than one card is produced from one source; the per-claim tag is carried by claim\_key.\par\medskip
\noindent\textit{Editorial note.} The author writes '(details omitted)' once, at proof line 37861, for the reduction to a finite free ambient module; that phrase is kept verbatim. The two references inside the counted proof span are recorded as non-load-bearing navigation pointers.\par\medskip
\noindent\textit{Allowed standard prerequisites (editorial).}\begin{itemize}
\item Bezout domains: every finitely generated ideal is principal. the hypothesis at statement line 37848 and used at proof lines 37862-37863. Is the input for the n = 1 case and for the rank-one pieces of the induction.
\item The Warfield decomposition for finitely presented modules over a ring whose localizations are valuation rings. invoked at proof lines 37873-37876. Produces the ambient module R\textasciicircum{}\{oplus r\} oplus bigoplus R/f\_iR of which M is a summand.
\item Functoriality of the torsion submodule and splitting of a direct summand sequence. used at proof lines 37877-37883. Separates the free and torsion parts and assembles the direct sum decomposition of M.
\end{itemize}\medskip
\ifdefined\hypertarget\hypertarget{sourcepointer-1}{}\fi
\noindent\textit{Source reference.} Original lemma 37834: sources/stacks-more-algebra/pinned-original-source.tex, source lines 37834--37834. Same-chapter navigation pointer; the statement it names is standard for the section and the argument that consumes it stays inside the counted proof span.\par\medskip
\ifdefined\hypertarget\hypertarget{sourcepointer-2}{}\fi
\noindent\textit{Source reference.} Original lemma 37737: sources/stacks-more-algebra/pinned-original-source.tex, source lines 37737--37737. Same-chapter navigation pointer; the statement it names is standard for the section and the argument that consumes it stays inside the counted proof span.\par\medskip
\begin{lemma}
\label{lemma-split-off-free-part}
Let $R$ be a B\'ezout domain.
\begin{enumerate}
\item Every finite submodule of a free module is finite free.
\item Every finitely presented $R$-module $M$ is a direct sum of a
finite free module and a torsion module $M_{tors}$ which is a
summand of a module of the form $\bigoplus_{i = 1, \ldots, n} R/f_iR$
with $f_1, \ldots, f_n \in R$ nonzero.
\end{enumerate}
\end{lemma}

\begin{proof}
Proof of (1). Let $M \subset F$ be a finite submodule of a free module $F$.
Since $M$ is finite, we may assume $F$ is a finite free module
(details omitted). Say $F = R^{\oplus n}$. We argue by induction
on $n$. If $n = 1$, then $M$ is a finitely generated ideal, hence
principal by our assumption that $R$ is B\'ezout. If $n > 1$, then
we consider the image $I$ of $M$ under the projection
$R^{\oplus n} \to R$ onto the last summand. If $I = (0)$, then
$M \subset R^{\oplus n - 1}$ and we are done by induction.
If $I \not = 0$, then $I = (f) \cong R$. Hence
$M \cong R \oplus \Ker(M \to I)$ and we are done by induction as well.

\medskip\noindent
Let $M$ be a finitely presented $R$-module. Since the localizations
of $R$ are maximal ideals are valuation rings
(Lemma \ref{lemma-localize-bezout})
we may apply Lemma \ref{lemma-warfield}.
Thus $M$ is a summand of a module of the form
$R^{\oplus r} \oplus \bigoplus_{i = 1, \ldots, n} R/f_iR$
with $f_i \not = 0$. Since taking the torsion submodule is
a functor we see that $M_{tors}$ is
a summand of the module $\bigoplus_{i = 1, \ldots, n} R/f_iR$
and $M/M_{tors}$ is a summand of $R^{\oplus r}$.
By the first part of the proof we see that $M/M_{tors}$ is
finite free. Hence $M \cong M_{tors} \oplus M/M_{tors}$
as desired.
\end{proof}

\end{document}
